% Avsnitt 12.1, ur lectures/L12/appendix/a_meta_prev.md.

\section{Att konstruera \code{meta_prev}}
\label{c12:sec:metaprev}

\code{rx_bus} anländer från en transceiver på ett annat chip, och \code{reset_n} från en knapp eller
en övervakare, så ingenting i den här FPGA:n har någonsin styrt någondera. Att sampla en av dem
direkt in i \code{clock}s domän kan lämna en vippa metastabil, vilket är problemet \chapref{c4:ch}
löste med två vippor i serie. Du har byggt den här modulen förut, som \code{sync} i
Övning~\ref{c5:ex:sync}; den här är samma idé, skriven en gång, ordentligt, för att behållas i den
här designen.

\Secref{c11:sec:sync} avgjorde redan \emph{att} kontrollern synkroniserar sina egna två asynkrona
ingångar, och varför: de är asynkrona vem som än instansierar \code{can_controller}, så kontrollern
äger dem och en integratör kopplar båda rakt till pinnar utan att tänka mer på det. Det här avsnittet
bygger modulen som gör det. Den är den enda modulen i hela boken som inte innehåller någonting alls
om CAN, vilket är precis därför den är återanvändbar nog att sluta med tre instanser vid
\chapref{c19:ch}.

\subsection{Gränssnitt}
\label{c12:sub:mpiface}

\modulefig[0.62]{lectures/L12/appendix/images/meta_prev.png}

\begin{booktable}{@{}p{5em}p{4.4em}p{11.4em}L@{}}
  \thead{Port / generic} & \thead{Riktn.} & \thead{Typ} & \thead{Betydelse} \\
  \midrule
  \code{WIDTH} & generic & \code{natural := 1} & Antal oberoende bitar att synkronisera. \\
  \code{clock} & in & \code{std_logic} & 50 MHz systemklocka, domänen som korsas \emph{in i}. \\
  \code{async_in} & in & \code{std_logic_vector(WIDTH-1 downto 0)} & Den eller de asynkrona
    ingångarna. \\
  \code{sync_out} & out & \code{std_logic_vector(WIDTH-1 downto 0)} & Samma bitar, säkra att använda
    i den här domänen. \\
\end{booktable}

\code{meta_prev_tb} binder sina portar positionellt, som allt i den här boken gör, så ordningen och
typerna ovan måste stämma exakt. \textbf{Generic-mappningen är också positionell}, så \code{WIDTH}
måste vara modulens enda generic (eller åtminstone dess första). Testbänken instansierar modulen två
gånger, en gång med standardbredden och en gång med bredd 4, så att både standardvärdet och genericen
övas. Den är också den enda modulen här som inte läser något paket alls, så den behöver varken
\code{can_def} eller något annat.

\subsection{Beteende}
\label{c12:sub:mpbehav}

Två vippor per bit, i serie, på \code{clock}s stigande flank: den första kan bli metastabil, och den
andra ger den en hel klockperiod att stabilisera sig innan något nedströms ser den.
\code{sync_out} följer därför \code{async_in} efter exakt två stigande flanker.

Tre saker skiljer sig från modulen \code{sync} du skrev i Övning~\ref{c5:ex:sync}, och alla tre är
avsiktliga:
\begin{itemize}
  \item \textbf{Ingen resetport.} \code{sync} tog en \code{reset_s2_n}. Den här tar ingen: efter två
    klockflanker håller en synkroniserare riktig data oavsett vad den startade från. Att stryka
    porten är också det som låter \emph{samma} modul synkronisera en reset, utan cirkulariteten i att
    en resetsynkroniserare behöver en synkroniserad reset. Testbänken klockar igenom den
    ursprungliga nivån innan den kontrollerar någonting, av precis det skälet.
  \item \textbf{Ingen generic \code{PRESET}.} \Chapref{c5:ch} argumenterade utförligt för att det
    säkra resetvärdet beror på vad signalen \emph{betyder}, och att bara den instansierande designen
    vet det. Det argumentet handlade om de två cyklerna efter reset, och det överlever inte att
    reseten stryks: utan resetgren finns inget värde att förinställa. Det som ersätter det är
    anroparens skyldighet att ignorera \code{sync_out} under de två första flankerna, vilket är samma
    skyldighet uttryckt annorlunda.
  \item \textbf{\code{SIZE} heter \code{WIDTH}.} Samma generic, samma uppgift, omdöpt till namnet
    varje delblock i det här projektet använder.
\end{itemize}

Den här designen instansierar den tre gånger, alla med \code{WIDTH = 1}: två gånger inuti
\code{can_controller}, för \code{reset_n} och för \code{rx_bus}, och en gång på \code{can_spi_node}s
toppnivå för dess egen reset (\chapref{c19:ch}). En modul, tre instanser, inga kopior. Genericen gör
ändå skäl för sin plats \textemdash{} testbänken övar \code{WIDTH = 4}, och en synkroniserare skriven
för en bit är en synkroniserare du kopierar nästa gång du behöver fyra.

\begin{warning}
En reservation att nämna redan nu, eftersom \chapref{c19:ch} lutar sig mot den: en synkroniserare
hanterar metastabilitet, inte \emph{skevhet}. Den andra vippan gör det överväldigande osannolikt att
ett metastabilt värde fortplantar sig (aldrig omöjligt; sannolikheten faller exponentiellt med
stabiliseringstiden, vilket är därför två vippor är den accepterade avvägningen), men varje bit
korsar fortfarande oberoende, så två bitar som ändras tillsammans kan landa en klocka isär. Det är
skälet till att en synkroniserare inte är ett generellt svar för en buss vars bitar måste hålla ihop
\textemdash{} och det är därför \code{can_spi_node} synkroniserar sin reset med den här modulen men
lämnar SPI-ledningarna till \code{spi_slave}, som ramar in dem, snarare än att bredda en
\code{meta_prev} över dem.
\end{warning}

\subsection{Vad testbänken låser fast}
\label{c12:sub:mptb}

\begin{itemize}
  \item \code{sync_out} följer \code{async_in} efter exakt två stigande flanker: inte en (en enda
    vippa är ingen synkroniserare) och inte tre.
  \item Detsamma gäller från hög till låg som från låg till hög.
  \item Med \code{WIDTH = 4} anländer alla fyra bitarna tillsammans, två flanker efter att ingången
    ändrats.
  \item En stabil ingång lämnar \code{sync_out} stabil.
\end{itemize}

Det här är den första testbänken i projektet som faktiskt \emph{körs} snarare än bara analyseras, så
det är också här det tredje GHDL-steget anländer. \Chapref{c11:ch} analyserade och elaborerade; se
bilaga~\ref{app:sim} för vad \code{ghdl -r} tillför och varför \code{--assert-level=error} inte är
valfritt.

\subsection{Att lägga in den i \code{can_controller}}
\label{c12:sub:mpwire}

\code{meta_prev}s portar är vektorer även vid bredd 1, så en skalär port behöver en enelementssignal
på vardera sidan. I arkitekturens deklarativa del (mellan \code{is} och \code{begin}):

\begin{vhdlcode}
    -- meta_prev, synchronizing the two inputs that arrive from off-chip.
    signal reset_n_v, reset_s2_n_v : std_logic_vector(0 downto 0);
    signal rx_bus_v,  rx_bus_s2_v  : std_logic_vector(0 downto 0);
    signal reset_s2_n, rx_bus_s2   : std_logic;
\end{vhdlcode}

och efter \code{begin}:

\begin{vhdlcode}
    reset_n_v(0) <= reset_n;
    rx_bus_v(0)  <= rx_bus;

    reset_s2_n <= reset_s2_n_v(0);
    rx_bus_s2  <= rx_bus_s2_v(0);

    reset_sync: entity work.meta_prev
        port map(clock, reset_n_v, reset_s2_n_v);

    rx_bus_sync: entity work.meta_prev
        port map(clock, rx_bus_v, rx_bus_s2_v);
\end{vhdlcode}

Två instanser i stället för en med \code{WIDTH => 2}. Båda bitarna skulle åka med tillsammans
alldeles utmärkt, eftersom de är oberoende och skevhet mellan dem inte betyder något, men en namngiven
instans per angelägenhet läser sig bättre på diagrammet, och att senare låta reseten hävdas asynkront
skulle då inte riva upp någonting.

Från och med nu \textbf{tar varje block \code{reset_s2_n}, aldrig \code{reset_n}, och varje block som
läser bussen läser \code{rx_bus_s2}, aldrig \code{rx_bus}}. Det är hela skälet till att de här två
instanserna finns, det gäller \code{bit_timer} redan i nästa avsnitt, och det är lätt att glömma igen
när \code{rx_shift_reg} kopplas in i \chapref{c15:ch}.
